\chapter{\texorpdfstring{Cross-Grid Input Compatibility and Scaling}{Cross-Grid Input Compatibility and Scaling}}
\label{ch:encoder-inputs}

A GNN can process grids of different sizes, but this alone does
not guarantee that its learned representation will transfer. The same input
feature must also have a comparable numerical meaning on both grids. For
example, an encoder trained on \texttt{bus14} may behave poorly on WCCI if a
power or angle feature has a different scale or reference. This chapter checks
that compatibility before evaluating transfer.
\\
\\
The comparison concerns the graph encoder used by the actors.
The critic already receives a normalized flat observation and remains specific
to each grid. By contrast, the transferred encoder initially reads raw graph
features because physical scaling and running standardization are disabled.
Its inputs must therefore be checked directly.


\section{Measurement protocol}
\label{sec:wcci-protocol}

A do-nothing policy is run for 6{,}000 environment steps on
each grid. Each chronic contributes at most 300 steps, which spreads the sample
over roughly twenty scenarios instead of concentrating it in one. At every
step, the local graph of each agent is rebuilt and all visible nodes are
recorded. All the agents are pooled because they use the same encoder. The final
sample contains 426{,}000 node observations for \texttt{bus14} and
1{,}236{,}000 for WCCI.
\\
\\
Only node features are compared. In this graph schema, the nine
edge channels only identify relation types and contain no continuous physical
measurements.
\\
\\
Three statistics describe the difference between each pair of
feature distributions. Let $A$ denote \texttt{bus14} and $B$ denote WCCI, with
empirical distributions $F_A$ and $F_B$, means $\mu_A$ and $\mu_B$, and
standard deviations $\sigma_A$ and $\sigma_B$. The pooled spread is
$\sigma_p=\sqrt{(\sigma_A^2+\sigma_B^2)/2}$.
\begin{align}
  \text{offset} &= \frac{\mu_B-\mu_A}{\sigma_p},
  &\text{scale ratio} &= \frac{\sigma_B}{\sigma_A},
  &\mathrm{KS} &= \sup_x |F_A(x)-F_B(x)| . \label{eq:offset-scale}
\end{align}
A positive offset means that WCCI has the higher mean. A scale
ratio above one means that WCCI has the wider distribution. The
Kolmogorov--Smirnov (KS) statistic is the largest difference between the two
empirical cumulative distributions and lies between zero and one.
\\
\\
These statistics describe different problems. The offset
detects a shift in the centre, while the scale ratio detects a change in spread.
KS summarizes the remaining disagreement. A feature is easiest to transfer
when its offset is near zero and its scale ratio is near one.


\section{What the encoder reads}
\label{sec:wcci-features}

The transfer study uses the disaggregated line-node graph. Its
twelve node channels are listed in
Table~\ref{tab:heterogeneous-line-node-features}. All node types use the same
feature vector, although most channels apply to only one type. The remaining
entries are set to zero, so the feature matrix is sparse by construction. For
example, \texttt{rho} is defined on transmission-line nodes and is zero on
busbars, generators and loads. These zeros indicate that the feature does not
apply to the node; they are not physical measurements equal to zero.
\\
\\
This sparsity can hide differences between the grids because
the shared zeros dominate the complete distribution. The analysis therefore
reports each feature both over all rows and over only the rows on which it is
defined. For example, the \texttt{p} and \texttt{theta} results combine
generator and load nodes, while \texttt{rho} uses transmission-line nodes.
Channels that remain constant under the do-nothing policy are kept in the table
for completeness, but they provide little information about cross-grid
compatibility.


\section{Results}
\label{sec:wcci-results}

Table~\ref{tab:wcci-shift} reports each statistic in two ways.
The first uses all sampled rows, including the structural zeros. The second
keeps only rows on which the feature is defined. The second view gives the
clearer comparison of the physical measurements.

\begin{table}[H]
  \centering
  \small
  \caption{Distribution differences between \texttt{bus14} and WCCI on the
  disaggregated line-node graph. Offset is measured in pooled standard
  deviations, and the scale ratio is WCCI divided by \texttt{bus14}. The final
  three columns keep only rows on which each measurement is defined. Type
  indicators and \texttt{domain\_mask} have no such restricted view. A constant
  channel has no defined scale ratio.}
  \label{tab:wcci-shift}
  \setlength{\tabcolsep}{4pt}
  \begin{tabular}{lrrrrrrrr}
    \toprule
    & \multicolumn{2}{c}{Zeros (\%)} & \multicolumn{3}{c}{All sampled rows}
    & \multicolumn{3}{c}{Applicable rows} \\
    \cmidrule(lr){2-3}\cmidrule(lr){4-6}\cmidrule(lr){7-9}
    Channel & \texttt{bus14} & WCCI & offset & ratio & KS & offset & ratio & KS \\
    \midrule
    \texttt{theta}  & 77.5 & 73.8 & $+0.653$ & 0.83 & 0.157 & $+1.864$ & 1.43 & \textbf{0.678} \\
    \texttt{p}      & 77.8 & 76.7 & $-0.037$ & 2.14 & 0.042 & $-0.158$ & \textbf{2.54} & 0.201 \\
    \texttt{rho}    & 63.4 & 63.6 & $-0.088$ & 0.92 & 0.053 & $-0.260$ & 1.05 & 0.141 \\
    \texttt{line\_status} & 63.4 & 63.6 & $-0.005$ & 1.00 & 0.002 & \multicolumn{3}{c}{constant at one} \\
    \texttt{timestep\_overflow} & 100 & 100 & $-0.002$ & 0.70 & 0.000 & $-0.004$ & 0.70 & 0.000 \\
    \texttt{time\_before\_cooldown\_line} & 100 & 100 & N/A & N/A & 0.000 & \multicolumn{3}{c}{constant} \\
    \texttt{time\_before\_cooldown\_sub}  & 100 & 100 & N/A & N/A & 0.000 & \multicolumn{3}{c}{constant} \\
    \midrule
    \texttt{is\_busbar}   & 60.6 & 65.0 & $-0.093$ & 0.98 & 0.045 & \multicolumn{3}{c}{N/A} \\
    \texttt{is\_generator} & 91.5 & 89.3 & $+0.076$ & 1.11 & 0.022 & \multicolumn{3}{c}{N/A} \\
    \texttt{is\_load}      & 84.5 & 82.0 & $+0.066$ & 1.06 & 0.025 & \multicolumn{3}{c}{N/A} \\
    \texttt{is\_transmission\_line} & 63.4 & 63.6 & $-0.004$ & 1.00 & 0.002 & \multicolumn{3}{c}{N/A} \\
    \texttt{domain\_mask}  & 0 & 0 & N/A & N/A & 0.000 & \multicolumn{3}{c}{constant at one} \\
    \bottomrule
  \end{tabular}
\end{table}

\noindent Over all rows, no channel reaches $\mathrm{KS}=0.2$ because the
many zeros hide the differences. On the applicable rows, \texttt{theta} has
$\mathrm{KS}=0.678$ and an offset of $+1.86$ standard deviations, while
\texttt{p} is $2.54$ times wider on WCCI. The type indicators differ little,
showing a similar mix of busbars, assets and transmission lines.
\\
\\
The raw ranges tell the same story. Power spans $[0,97]$~MW on
\texttt{bus14} and $[-281,399]$~MW on WCCI, where $7.8\,\%$ of asset rows are
negative. These values come mainly from WCCI load profiles: a negative load is
a net injection rather than consumption. They are retained because the sign is
physical, and they lie outside the range seen during \texttt{bus14} training.
Angles span $[-11.6,1.5]$ and $[-7.1,16.8]$ degrees. Power therefore differs in
both scale and support, while the angle distribution mainly shifts its centre.

\begin{figure}[H]
  \centering
  \includegraphics[width=\linewidth]{figures/wcci_feature_distributions.png}
  \caption{Feature distributions on the rows where each channel is defined. \texttt{rho} is
  similar on both grids. \texttt{p} has a much wider range on WCCI, while
  \texttt{theta} is mainly shifted rather than widened. The last panel reports
  the KS statistic.}
  \label{fig:wcci-distributions}
\end{figure}


\section{Scale changes and offsets}
\label{sec:wcci-interpretation}

Figure~\ref{fig:wcci-offset-scale} separates a change in centre
from a change in spread. Each point represents one feature. The horizontal axis
shows the absolute offset, and the vertical axis shows the scale ratio on a
logarithmic scale. A feature near the bottom-left corner has a similar centre
and spread on both grids. The left panel shows the raw inputs, while the right
panel shows the transformed inputs discussed in
Section~\ref{sec:wcci-corrected}.
\\
\\

\begin{figure}[htbp]
  \centering
  \includegraphics[width=\linewidth]{figures/wcci_feature_offset_scale.png}
  \caption{Absolute offset against scale ratio on the rows where each channel
  is defined. The
  shaded band marks scale ratios close to one, and the dotted vertical line
  marks an offset of $0.3$ pooled standard deviations. The left panel uses raw
  features. The right panel uses the transformations from
  Section~\ref{sec:wcci-correction}. The largest offset falls from $1.86$ to
  $0.81$, although \texttt{p} gains a larger offset.}
  \label{fig:wcci-offset-scale}
\end{figure}

\newpage

\noindent The raw distributions reveal two different
mismatches. The \texttt{p} channel differs in scale and includes negative WCCI
values that were not seen during source training, while \texttt{theta} has an
offset of $1.86$ pooled standard deviations. These observations motivate
testing deterministic power scaling and a reference-independent angle
representation.
\\
\\
The \texttt{rho} channel is already comparable. Its scale ratio
is $1.05$ and its offset is $-0.26$. This is expected because \texttt{rho}
divides the loading of each transmission line by its own thermal limit. It shows
the value of using grid-relative quantities when the physics allows it.


\section{Transforming the inputs}
\label{sec:wcci-correction}

We explore provides two forms of preprocessing. Physical scaling
divides a feature by a fixed constant defined from grid metadata. Running
standardization instead estimates a mean and standard deviation during
training. The implementation keeps separate statistics for each node type, so
the sparse feature blocks do not contaminate one another. It standardizes
continuous channels such as \texttt{p} and \texttt{theta}, but deliberately
leaves \texttt{rho} unchanged.

\subsection{Running standardization does not align a new grid}
\label{sec:wcci-running-norm}

During source training, running standardization can centre and
scale the \texttt{bus14} inputs and the resulting statistics are saved. In strict zero-shot transfer, they must remain frozen: both grids
use the same transformation $(x-\mu_{14})/\sigma_{14}$. This centres
\texttt{bus14}, but WCCI keeps its relative offset and scale difference.
Without clipping, applying the same monotonic affine transformation to both
grids also leaves their KS distance unchanged. Running standardization may help
source training, but it is not a cross-grid alignment method.

\subsection{Physical scaling for the power channels}
\label{sec:wcci-physical-scaling}

Physical scaling divides the power channels by the largest
generator capacity in each environment, as defined in
Section~\ref{subsec:graph-feature-normalization}. Scaling leaves the
structural zeros at zero, and the divisor is fixed before training. The largest
generator is only an approximate measure of grid size, so
Table~\ref{tab:wcci-normalizer} compares it with several alternatives. Among
the tested choices, it brings the power spreads closest to equality.

\begin{table}[H]
  \centering
  \small
  \caption{Candidate power normalizers and the resulting \texttt{p} scale ratio
  on asset nodes. A ratio of $1.00$ is ideal. Without scaling, the ratio is
  $2.54$.}
  \label{tab:wcci-normalizer}
  \setlength{\tabcolsep}{8pt}
  \begin{tabular}{lrrr}
    \toprule
    Divisor & \texttt{bus14} & WCCI & Resulting ratio \\
    \midrule
    \textbf{Maximum $P_{\max}$ (used)} & 140 & 400 & \textbf{0.89} \\
    Total $P_{\max}$ & 540 & 2{,}450 & 0.56 \\
    Generator count & 6 & 22 & 0.69 \\
    Median $P_{\max}$ & 85 & 71 & 3.04 \\
    \bottomrule
  \end{tabular}
\end{table}

\subsection{Replacing absolute voltage angles}
\label{sec:wcci-angle-representation}

Scaling cannot solve the angle problem. Absolute voltage angles
are measured relative to the slack bus, and the two grids use different slack
locations. The shift in Figure~\ref{fig:wcci-distributions} is therefore partly
caused by the reference choice rather than by a physical difference between the
grids.
\\
\\
The solution is to use the angle difference across each
transmission line, as described in Section~\ref{sec:line-node-angle}. This
difference does not depend on the slack-bus reference. It is stored on the
transmission-line node, which is also more common than the generator and load
nodes that carry the absolute angle. Both environments therefore replace
absolute angles with line-angle differences. This changes both the physical
quantity and the node type that carries it; it is a new representation rather
than a rescaling of the original channel.

\subsection{Results of the transformations}
\label{sec:wcci-corrected}

Table~\ref{tab:wcci-corrected} evaluates both transformations
under the same protocol as the raw inputs. The right panel of
Figure~\ref{fig:wcci-offset-scale} shows the result. The angle difference now
behaves similarly to \texttt{rho}, and the power-scale ratio is within
$11\,\%$ of one.

\begin{table}[H]
  \centering
  \small
  \caption{Distribution differences for the original and transformed inputs
  after enabling
  \texttt{gnn\_angle\_representation = edge\_diff} and
  \texttt{gnn\_physical\_scaling = true}. Statistics are computed on the node
  rows where each channel is defined. The offset is retained because scaling
  each grid by a different constant can also move the centre.}
  \label{tab:wcci-corrected}
  \setlength{\tabcolsep}{6pt}
  \begin{tabular}{lrrr|rrr}
    \toprule
    & \multicolumn{3}{c|}{Original} & \multicolumn{3}{c}{Transformed} \\
    \cmidrule(lr){2-4}\cmidrule(lr){5-7}
    Channel & offset & ratio & KS & offset & ratio & KS \\
    \midrule
    \texttt{theta}      & $+1.864$ & 1.43 & 0.678 & \multicolumn{3}{c}{replaced} \\
    \texttt{theta\_diff} & \multicolumn{3}{c|}{N/A} & $-0.383$ & \textbf{1.15} & \textbf{0.230} \\
    \texttt{p}          & $-0.158$ & 2.54 & 0.201 & $-0.808$ & \textbf{0.89} & 0.511 \\
    \texttt{rho}        & $-0.260$ & 1.05 & 0.141 & $-0.260$ & 1.05 & 0.141 \\
    \bottomrule
  \end{tabular}
\end{table}

\noindent The angle transformation clearly improves cross-grid
compatibility. Replacing the absolute angle with a reference-independent
difference reduces KS from $0.678$ to $0.230$, while the scale ratio remains
close to one. Because the two channels describe different quantities on
different node types, this is a comparison between representations rather than
a direct normalization of the same variable.
\\
\\
The power transformation involves a trade-off. Scaling
improves the spread ratio from $2.54$ to $0.89$, but it increases the offset from $-0.158$
to $-0.808$ and raises KS from $0.201$ to $0.511$. The reason is that
the mean sampled generator output represents about $39\,\%$ of the largest
generator capacity on \texttt{bus14}, but only $20\,\%$ on WCCI. These values
do not describe the utilization of each individual generator. Dividing by the
largest generator capacity therefore places typical outputs at different
numerical levels on the two grids.
\\
\\
The evidence supports the angle representation, while power
scaling remains uncertain. The transformed-input condition combines both
transformations as an experimental condition, while the raw-input condition
keeps the original features as a control.

\begin{figure}[H]
  \centering
  \includegraphics[width=\linewidth]{figures/wcci_feature_corrected.png}
  \caption{Transformed feature distributions. The angle difference is centred in
  a similar region on both grids, and the scaled power channels now have similar
  widths. Compare with the raw distributions in
  Figure~\ref{fig:wcci-distributions}.}
  \label{fig:wcci-corrected}
\end{figure}

\section{Inputs retained for transfer}
\label{sec:wcci-corrected-inputs}

The transfer experiments retain two input conditions. The raw-input condition
uses the original graph features. The transformed-input condition divides power
by the largest generator on each grid and replaces absolute voltage angles with
line-angle differences. Running standardization is disabled because frozen
source statistics cannot align WCCI in strict zero-shot transfer. Comparing
otherwise matched raw-input and transformed-input models isolates the effect of
deterministic input preprocessing. The later bridge experiments separately
study the effect of source pretraining in
Chapters~\ref{ch:encoder-transfer-bridge} and~\ref{ch:screen-f-transfer}.
\\
\\
\begin{tcolorbox}[
  colframe=red!75!black,
  colback=red!5!white,
  title=Chapter 8 Summary: Cross-Grid Input Compatibility and Scaling]
\textbf{Purpose:} Check whether the shared encoder receives
comparable numerical inputs on \texttt{bus14} and WCCI.
\\
\\
\textbf{Content:} Do-nothing rollouts provide $426{,}000$
\texttt{bus14} and $1{,}236{,}000$ WCCI graph-node entries. Each feature is
compared using its offset, scale ratio and KS statistic.
\\
\\
\textbf{Findings:}
\begin{itemize}
  \item \texttt{rho} is already comparable because each line
  loading is divided by its own thermal limit.
  \item WCCI power has a wider range and includes negative net
  injections that are absent from the sampled \texttt{bus14} inputs.
  \item Replacing absolute angles with line-angle differences
  reduces angle KS from $0.678$ to $0.230$.
  \item Power scaling improves the spread ratio from $2.54$ to
  $0.89$, but increases the offset from $-0.158$ to $-0.808$.
  \item Type-aware running standardization preserves sparse
  zeros, but frozen source statistics cannot align WCCI in strict zero-shot
  transfer.
\end{itemize}
\textbf{Takeaway:} The angle transformation is supported by the
distribution analysis, whereas power scaling has mixed effects. The raw-input
and transformed-input conditions are therefore both evaluated. Target-specific
standardization would be a separate calibration experiment rather than strict
zero-shot transfer.
\end{tcolorbox}
